\chapterhead{43}{ORR-4}{Risk Measurement and Assessment}{1}{2}

\lohead{43}{a}{Explain best practices for the collection of operational loss data
and reporting of operational loss incidents, including regulatory expectations.}

\orsub{Regulatory requirements for operational loss data}

\noindent\begin{minipage}{\textwidth}
\renewcommand{\arraystretch}{1.3}
\begin{tabularx}{\textwidth}{@{}L{0.31\textwidth}Y@{}}
\rowcolor{navy}\multicolumn{2}{@{}l@{}}{\hspace{4pt}\thd{Minimum standards}}\\
{\tabfont \tsub{Minimum data collection period}} & {\tabfont At least 10 years}\\
\rowcolor{paleblue}
{\tabfont \tsub{Minimum collection threshold}} & {\tabfont EUR 20,000 (many banks use a lower threshold)}\\
{\tabfont \tsub{Loss classification}} & {\tabfont Based on Basel event types}\\
\rowcolor{paleblue}
{\tabfont \tsub{Reported dates}} & {\tabfont Occurrence, discovery, reporting, and accounting}\\
\rowcolor{navy}\multicolumn{2}{@{}l@{}}{\hspace{4pt}\thd{Data collection process --- how data should be collected}}\\
{\tabfont \tsub{Documented procedures}} & {\tabfont For identification, collection, and treatment of data}\\
\rowcolor{paleblue}
{\tabfont \tsub{Ensure clear reporting}} & {\tabfont Dates, types, thresholds, and quantification}\\
{\tabfont \tsub{Standardised approach}} & {\tabfont Consistent taxonomy for gathering information}\\
\rowcolor{navy}\multicolumn{2}{@{}l@{}}{\hspace{4pt}\thd{Comprehensive data --- cover all data possible}}\\
{\tabfont \tsub{Coverage}} & {\tabfont Collecting data from all significant activities and locations}\\
\rowcolor{paleblue}
{\tabfont \tsub{Loss threshold}} & {\tabfont Many banks use a lower threshold than EUR 20,000, for comprehensiveness}\\
{\tabfont \tsub{Cost inclusion}} & {\tabfont Include both direct (financial loss) and indirect costs (reputation damage)}\\
\rowcolor{paleblue}
{\tabfont \tsub{Loss grouping}} & {\tabfont Group multiple losses stemming from a common cause}\\
\rowcolor{navy}\multicolumn{2}{@{}l@{}}{\hspace{4pt}\thd{Incident dates --- reporting operational loss details}}\\
{\tabfont \tsub{Reported information}} & {\tabfont Report gross loss, net recoveries, and reference dates}\\
\rowcolor{paleblue}
{\tabfont \tsub{Mention relevant dates}} & {\tabfont Occurrence (when it first happened), discovery, reporting, and accounting}\\
{\tabfont \tsub{Time interval tracking}} & {\tabfont For internal analysis. A huge interval between occurrence and discovery reflects an issue with visibility}\\
\rowcolor{navy}\multicolumn{2}{@{}l@{}}{\hspace{4pt}\thd{Boundary events --- market/credit or operational loss?}}\\
{\tabfont \tsub{Classification basis}} & {\tabfont These events are classified based on \textbf{cause}, not consequence}\\
\rowcolor{paleblue}
{\tabfont \tsub{Credit losses due to operational errors}} & {\tabfont Included in RWAs}\\
{\tabfont \tsub{Market losses due to operational errors}} & {\tabfont Included in operational losses}\\
\end{tabularx}
\end{minipage}
\orfigcap{Regulatory expectations for operational loss data collection.}

\orsub{Data quality requirements}

\noindent\begin{minipage}{\textwidth}
\renewcommand{\arraystretch}{1.35}
\begin{tabularx}{\textwidth}{@{}L{0.30\textwidth}Y@{}}
\rowcolor{navy} \thd{Requirement} & \thd{What it means}\\
{\tabfont \textbf{Data completeness and accuracy}} & {\tabfont Must have processes to ensure the completeness and accuracy of data}\\
\rowcolor{paleblue}
{\tabfont \textbf{Corroboration of information}} & {\tabfont Use multiple data sources, including the general ledger and IT logs}\\
{\tabfont \textbf{Mitigate underreporting}} & {\tabfont Implement strategies like reminders, performance metrics, and internal audits}\\
\end{tabularx}
\end{minipage}
\orfigcap{Three data quality requirements.}

\orsub{Challenges of operational risk data}

\begin{lettered}
\item Operational risk data is idiosyncratic and less correlated to markets
compared to credit and market risk data.
\item It has a wide distribution of returns with more frequent extreme events.
\item The interrelated nature of operational risk makes cause-and-effect analysis
difficult.
\end{lettered}

% ---------------------------------------------------------------------
\lohead{43}{b}{Explain operational risk-assessment processes and tools, including
risk control self-assessments (RCSAs), likelihood assessment scales, and
heatmaps.}

\orsub{Risk and control self-assessments}

\begin{orkeybox}
{\sffamily\bfseries\fontsize{8.6}{10}\selectfont\color{navy}Purpose}\par
\vspace{2pt}
To assess the probability and severity of operational risks for a bank. It aims
to increase the bank's awareness of inherent risks, internal controls, and
residual risks.\par
\vspace{4pt}
{\fontsize{9}{11}\selectfont\color{rust}\itshape Note: it does not quantify risk.}
\end{orkeybox}

\begin{checks}
\item \textbf{Inherent risk:} examines risks without considering existing
controls.
\item \textbf{Residual risk:} evaluates risks remaining after considering control
effectiveness.
\item \textbf{Frequency:} annually, and quarterly for very significant risks.
\end{checks}

\orsubb{RCSA process}

Participants discuss inherent risks, controls in place, and assess the controls'
overall effectiveness, sometimes using questionnaires. Two main types of controls:
\textbf{preventive} (reduce risk likelihood) and \textbf{corrective} (minimise
negative effects).

\orsubb{Limitations}

Subjectivity and behavioural biases, limited data, and comparison of results can
all be problematic.

\orsubb{RCSA output and mitigation}

\begin{itemize}
\item Documented results are sent to unit managers for approval and
implementation of mitigation actions.
\item RCSAs help determine if a unit's risk level aligns with the bank's overall
risk appetite.
\item Mitigation may involve new controls, control enhancements, reduced
transaction volumes, or insurance.
\end{itemize}

\orsub{Impact scales}

To assess the potential consequences of an operational risk event across various
categories like financial, regulatory, customer, and reputation. Scales are
typically 4--5 ratings (e.g., low, medium, high, very high, extreme).

\orsub{Likelihood assessment scales}

Likelihood assessment scales express the probability or frequency of risk
occurrence.

\noindent\begin{minipage}{\textwidth}
\renewcommand{\arraystretch}{1.35}
\begin{tabularx}{\textwidth}{@{}L{0.20\textwidth}YL{0.30\textwidth}@{}}
\rowcolor{navy} \thd{Qualitative rating} & \thd{Frequency of occurrence} & \thd{Probability of occurrence (one year horizon)}\\
\rowcolor{palered}
{\tabfont \textbf{Likely}} & {\tabfont Once per year or more frequently} & {\tabfont $>50\%$}\\
\rowcolor{palegold}
{\tabfont \textbf{Possible}} & {\tabfont $>1$--5 years} & {\tabfont 20--50\%}\\
\rowcolor{palegold}
{\tabfont \textbf{Unlikely}} & {\tabfont 1 in 5--20 years} & {\tabfont 5--20\%}\\
\rowcolor{palegreen}
{\tabfont \textbf{Remote}} & {\tabfont $<1$ in 20 years} & {\tabfont $<5\%$}\\
\end{tabularx}
\end{minipage}
\orfigcap{A four-point likelihood assessment scale.}

\orsub{Heatmap}

A heatmap, known as an RCSA matrix, combines likelihood and impact ratings to
visualise risk intensity using colours.

\begin{center}
\begin{tikzpicture}[font=\fontsize{8.2}{10}\selectfont,
  cell/.style={minimum width=2.35cm,minimum height=0.95cm,draw=white,line width=1.4pt},
  lab/.style={font=\sffamily\bfseries\fontsize{8}{9.6}\selectfont,text=navy}]

\definecolor{hgreen}{HTML}{57A773}
\definecolor{hyellow}{HTML}{F2D14E}
\definecolor{horange}{HTML}{ED9B54}
\definecolor{hred}{HTML}{D9534F}

\foreach \r/\row in {%
  0/{hyellow,horange,hred,hred},
  1/{hgreen,hyellow,horange,hred},
  2/{hgreen,hgreen,hyellow,horange},
  3/{hgreen,hgreen,hgreen,horange}}
{
  \foreach \c [count=\i from 0] in \row
    \node[cell,fill=\c] at (\i*2.4,-\r*1.0) {};
}
\foreach \r/\t in {0/Likely,1/Possible,2/Unlikely,3/Remote}
  \node[lab,anchor=east] at (-1.35,-\r*1.0) {\t};
\foreach \i/\t in {0/Low,1/Moderate,2/Major,3/Extreme}
  \node[lab] at (\i*2.4,-3.75) {\t};
\node[lab,anchor=east,rotate=90] at (-2.9,-1.5) {Likelihood};
\node[lab] at (3.6,-4.35) {Impact};
\end{tikzpicture}
\end{center}
\orfigcap{The RCSA matrix, or heatmap.}

% ---------------------------------------------------------------------
\lohead{43}{d}{Describe and distinguish between the different quantitative
approaches and models used to analyze operational risk.}

\lohead{43}{e}{Estimate operational risk exposures based on the fault tree model
given probability assumptions.}

\orsub{Quantitative analysis of operational risks}

\orsubb{1) Root-cause analysis}

\begin{lettered}
\item Conducted by the first line of defense, it aims to identify the root cause
of failures through in-depth investigation.
\item Repeatedly asks ``5 why'' to uncover multiple contributing causes.
\item Offers a thorough and methodical assessment of risk severity after
mitigation.
\end{lettered}

\begin{center}
\begin{tikzpicture}[font=\fontsize{7.4}{9}\selectfont,
  cz/.style={draw=palenavy,fill=paleblue,rounded corners=2pt,inner sep=3.5pt,
             text width=1.55cm,align=center,text=navy},
  ctl/.style={font=\sffamily\fontsize{6.8}{8.4}\selectfont,text=midgray}]

\node[circle,fill=navy,text=white,inner sep=4pt,align=center,
      font=\sffamily\bfseries\fontsize{7.6}{9}\selectfont] (ev) at (0,0) {Risk\\Event};

\node[cz] (dc1) at (-2.6, 0.85) {Direct Cause 1};
\node[cz] (dc2) at (-2.6,-0.85) {Direct Cause 2};
\node[cz] (ic1) at (-5.5, 1.35) {Indirect Cause 1};
\node[cz] (ic2) at (-5.5, 0.0)  {Indirect Cause 2};
\node[cz] (ic3) at (-5.5,-1.35) {Indirect Cause 3};

\node[cz] (di1) at ( 2.6, 0.85) {Direct Impact 1};
\node[cz] (di2) at ( 2.6,-0.85) {Direct Impact 2};
\node[cz] (ii1) at ( 5.5, 1.35) {Indirect Impact 1};
\node[cz] (ii2) at ( 5.5, 0.0)  {Indirect Impact 2};
\node[cz] (ii3) at ( 5.5,-1.35) {Indirect Impact 3};

\foreach \a/\b in {ic1/dc1,ic2/dc1,ic2/dc2,ic3/dc2}
  \draw[navy!45,line width=0.7pt] (\a) -- (\b);
\foreach \a in {dc1,dc2} \draw[-{Stealth[length=4pt]},navy,line width=0.9pt] (\a) -- (ev);
\foreach \a in {di1,di2} \draw[-{Stealth[length=4pt]},rust,line width=0.9pt] (ev) -- (\a);
\foreach \a/\b in {di1/ii1,di1/ii2,di2/ii2,di2/ii3}
  \draw[rust!45,line width=0.7pt] (\a) -- (\b);

\node[ctl] at (-4.0,-2.25) {Preventative controls};
\node[ctl] at ( 0.0,-2.25) {Detective controls};
\node[ctl] at ( 4.0,-2.25) {Corrective controls};
\draw[palenavy,line width=0.8pt] (-6.6,-1.95) -- (6.6,-1.95);
\end{tikzpicture}
\end{center}
\orfigcap{Root-cause, or bowtie, analysis: causes on the left, impacts on the
right, controls underneath.}

\orsubb{2) Swiss cheese model}

\begin{lettered}
\item Recognises that controls, like slices of cheese, may have holes
(weaknesses).
\item Requires independent control failures and compensating controls to mitigate
weaknesses.
\end{lettered}

\begin{center}
\begin{tikzpicture}[font=\fontsize{7.6}{9.4}\selectfont]
\foreach \i/\x in {1/0,2/1.9,3/3.8,4/5.7}{
  \begin{scope}[shift={(\x,0)}]
    \draw[fill=palegold,draw=gold,line width=0.8pt,rounded corners=2pt]
      (0,0) -- (1.1,0.55) -- (1.1,2.75) -- (0,2.2) -- cycle;
  \end{scope}}
\foreach \x/\y/\r in {0.50/1.28/0.17, 0.80/2.05/0.12,
                      2.42/1.32/0.15, 2.70/0.85/0.12,
                      4.32/1.36/0.17, 4.60/2.15/0.12,
                      6.22/1.40/0.15, 6.50/0.90/0.13}
  \fill[white] (\x,\y) circle (\r);
\draw[-{Stealth[length=5pt]},rust,line width=1.3pt]
  (-1.1,1.25) -- (7.7,1.25);
\node[text=rust,font=\sffamily\bfseries\fontsize{7.6}{9}\selectfont,anchor=east]
  at (-1.2,1.25) {Hazards};
\node[text=rust,font=\sffamily\bfseries\fontsize{7.6}{9}\selectfont,anchor=west]
  at (7.85,1.25) {Losses};
\node[text=navy,font=\sffamily\fontsize{7.2}{9}\selectfont] at (3.3,-0.55)
  {successive control layers, each with weaknesses};
\end{tikzpicture}
\end{center}
\orfigcap{The Swiss cheese model: a loss occurs when the holes line up.}

\orsubb{3) Fault tree analysis (FTA)}

\begin{lettered}
\item Breaks down failure scenarios into contributing factors (AND/OR
conditions).
\item Useful for understanding how a chain of failures can lead to a major loss.
\end{lettered}

\begin{center}
\begin{tikzpicture}[font=\fontsize{7.6}{9.4}\selectfont,
  gate/.style={diamond,fill=navy,text=white,inner sep=2.8pt,aspect=1.15,
               font=\sffamily\bfseries\fontsize{7.6}{9}\selectfont},
  esc/.style={text=midgray,font=\fontsize{7}{8.6}\selectfont,anchor=west},
  pl/.style={text=navy,font=\fontsize{7}{8.6}\selectfont,anchor=east}]

\foreach \i/\y/\lab/\p in {1/0/{Phishing email}/{P_1},
                           2/-1.25/{Firewall failure}/{P_2},
                           3/-2.5/{Click}/{P_3},
                           4/-3.75/{Activity detection}/{P_4},
                           5/-5.0/{Exit}/{P_5}}
{
  \node[gate] (g\i) at (0,\y) {\i};
  \node[pl] at (-0.65,\y) {\lab\ $\p$};
  \node[esc] at (1.05,\y) {No failure at step \i\ \ $(1-\p)$};
  \draw[-{Stealth[length=4pt]},midgray,line width=0.8pt] (g\i.east) -- ++(0.7,0);
}
\foreach \a/\b in {1/2,2/3,3/4,4/5}
  \draw[-{Stealth[length=4pt]},navy,line width=1pt] (g\a.south) -- (g\b.north);
\node[fill=navy,text=white,rounded corners=2.5pt,inner sep=5pt,
      font=\sffamily\bfseries\fontsize{8}{9.6}\selectfont] (inc) at (0,-6.25) {Incident};
\draw[-{Stealth[length=4pt]},navy,line width=1pt] (g5.south) -- (inc.north);
\end{tikzpicture}
\end{center}
\orfigcap{A fault tree: the incident occurs only when every stage fails in sequence.}

% ---------------------------------------------------------------------
\lohead{43}{f}{Describe approaches used to determine the level of operational risk
capital for economic capital purposes, including their application and
limitations.}

\orsub{Loss Distribution Approach (LDA)}

Breaks down loss events into \textbf{frequency} (how often) and \textbf{severity}
(how much), then uses simulations (Monte Carlo) to create a combined loss
distribution.

\begin{center}
\begin{tikzpicture}[font=\fontsize{7.2}{8.8}\selectfont,
  ttl/.style={font=\sffamily\bfseries\fontsize{7.6}{9}\selectfont,text=navy}]

% ---- frequency (left)
\begin{scope}[shift={(-4.6,0)}]
  \draw[midgray,line width=0.6pt] (0,0) -- (4.0,0);
  \foreach \x/\h in {0.4/0.35,0.8/0.75,1.2/1.15,1.6/1.35,2.0/1.05,2.4/0.7,2.8/0.4,3.2/0.2}
    \draw[navy,line width=2.6pt] (\x,0) -- (\x,\h);
  \node[ttl] at (2.0,2.45) {Loss frequency};
  \node[text=midgray,anchor=north] at (2.0,-0.12) {events per year};
\end{scope}

% ---- severity (right)
\begin{scope}[shift={(1.0,0)}]
  \draw[midgray,line width=0.6pt] (0,0) -- (4.0,0);
  \draw[teal,line width=1.4pt] plot[smooth,domain=0.1:3.9,samples=60]
    (\x, {2.0*exp(-((\x-0.75)^2)/0.28)+0.55*exp(-((\x-1.9)^2)/2.6)});
  \node[ttl] at (2.0,2.45) {Loss severity};
  \node[text=midgray,anchor=north] at (2.0,-0.12) {loss size (\$M)};
\end{scope}

% ---- arrows converging
\draw[-{Stealth[length=5pt]},navy!45,line width=1.2pt] (-2.6,-0.45) -- (-0.3,-1.35);
\draw[-{Stealth[length=5pt]},navy!45,line width=1.2pt] (3.0,-0.45) -- (0.7,-1.35);
\node[text=rust,font=\sffamily\bfseries\fontsize{7.4}{9}\selectfont] at (0.2,-1.0)
  {Monte Carlo};

% ---- aggregate
\begin{scope}[shift={(-2.0,-4.3)}]
  \draw[midgray,line width=0.6pt] (0,0) -- (4.4,0);
  \draw[midgray,line width=0.6pt] (0,0) -- (0,2.1);
  \foreach \x/\h in {0.3/0.5,0.6/1.1,0.9/1.7,1.2/1.95,1.5/1.6,1.8/1.15,
                     2.1/0.8,2.4/0.55,2.7/0.38,3.0/0.26,3.3/0.18,3.6/0.12,3.9/0.07}
    \draw[teal,line width=2.4pt] (\x,0) -- (\x,\h);
  \node[ttl] at (2.2,2.4) {Aggregate loss distribution};
  \node[text=midgray,anchor=north] at (2.2,-0.12) {loss (\$M)};
  \node[text=midgray,rotate=90,anchor=south] at (-0.3,1.05) {probability};
\end{scope}
\end{tikzpicture}
\end{center}
\orfigcap{Frequency and severity are convolved into a single aggregate loss
distribution.}

\begin{orexambox}{Strengths}
\begin{itemize}
\item Uses internal loss data, potentially more relevant to the bank.
\item Provides a comprehensive view of potential losses.
\end{itemize}
\end{orexambox}

\begin{ornotebox}{Weaknesses}
\begin{itemize}
\item Assumes independence between loss events, which may not be realistic.
\item Relies heavily on data quality and quantity.
\item Doesn't handle very high skewness (uneven distribution) of operational
losses well.
\item Difficulty in clustering data into homogeneous categories, units of
measurement (UoMs), with similar drivers.
\end{itemize}
\end{ornotebox}

\orsub{Extreme Value Theory (EVT)}

Focuses on analysing the ``tail'' of the loss distribution (very large losses).
Useful for identifying potential for very large losses.

\begin{center}
\begin{tikzpicture}[font=\fontsize{7.2}{8.8}\selectfont,
  ttl/.style={font=\sffamily\bfseries\fontsize{7.6}{9}\selectfont,text=navy},
  sm/.style={circle,draw=midgray,fill=white,inner sep=1.1pt},
  bg/.style={circle,fill=gold,inner sep=1.8pt}]

% ---- block maxima
\begin{scope}[shift={(-8.0,0)}]
  \draw[midgray,line width=0.6pt] (0,0) -- (6.2,0);
  \draw[midgray,line width=0.6pt] (0,0) -- (0,2.9);
  \foreach \b in {1.55,3.1,4.65}
    \draw[palenavy,line width=0.7pt,dashed] (\b,0) -- (\b,2.9);
  \foreach \x/\y in {0.3/0.35,0.55/0.7,0.8/0.3,1.05/0.55,1.3/0.25}
    \node[sm] at (\x,\y) {};
  \node[bg] at (0.7,1.75) {};
  \foreach \x/\y in {1.85/0.4,2.1/0.6,2.35/0.3,2.6/0.75,2.85/0.35}
    \node[sm] at (\x,\y) {};
  \node[bg] at (2.3,2.2) {};
  \foreach \x/\y in {3.4/0.3,3.65/0.65,3.9/0.4,4.15/0.3,4.4/0.6}
    \node[sm] at (\x,\y) {};
  \node[bg] at (3.85,1.5) {};
  \foreach \x/\y in {4.95/0.45,5.2/0.3,5.45/0.7,5.7/0.35,5.95/0.5}
    \node[sm] at (\x,\y) {};
  \node[bg] at (5.35,2.0) {};
  \node[ttl] at (3.1,3.25) {Block maxima (Fisher--Tippett)};
  \node[text=midgray,anchor=north] at (3.1,-0.12) {Time};
  \node[text=midgray,rotate=90,anchor=south] at (-0.35,1.45) {Loss amount};
\end{scope}

% ---- peaks over threshold
\begin{scope}[shift={(1.0,0)}]
  \draw[midgray,line width=0.6pt] (0,0) -- (6.2,0);
  \draw[midgray,line width=0.6pt] (0,0) -- (0,2.9);
  \draw[rust,line width=1pt] (0,1.25) -- (6.2,1.25);
  \node[text=rust,anchor=west,font=\sffamily\fontsize{7}{8.6}\selectfont]
    at (6.25,1.25) {$u$};
  \foreach \x/\y in {0.35/0.4,0.7/0.75,1.05/0.3,1.4/0.6,1.75/0.35,2.1/0.8,
                     2.45/0.45,2.8/0.3,3.15/0.7,3.5/0.4,3.85/0.55,4.2/0.3,
                     4.55/0.75,4.9/0.35,5.25/0.6,5.6/0.4,5.95/0.5}
    \node[sm] at (\x,\y) {};
  \foreach \x/\y in {0.6/1.85,1.5/2.3,2.4/1.6,3.3/2.05,4.3/1.75,5.4/2.4,5.9/1.55}
    \node[bg] at (\x,\y) {};
  \node[ttl] at (3.1,3.25) {POT (peaks-over-threshold)};
  \node[text=midgray,anchor=north] at (3.1,-0.12) {Time};
  \node[text=midgray,rotate=90,anchor=south] at (-0.35,1.45) {Loss amount};
\end{scope}
\end{tikzpicture}
\end{center}
\orfigcap{Two ways of isolating the tail: one maximum per block, or every
observation above a threshold.}

\begin{ornotebox}{Limitations}
\begin{itemize}
\item Requires significant data on large losses, which may not be available.
\item Assumes a single main cause for operational losses, often unrealistic.
\end{itemize}
\end{ornotebox}

\orsub{Internal and external loss data}

\begin{itemize}
\item Internal data provides detailed information on past incidents and controls.
\item External data from other banks broadens the understanding of potential loss
outcomes.
\end{itemize}

\begin{ornotebox}{Challenges}
\begin{checks}
\item Internal data may be incomplete; not all incidents are reported.
\item External data availability and comparability can be limited.
\item Combining data requires adjustments for factors like size and inflation.
\end{checks}
\end{ornotebox}

\orsub{Capital modeling}

\begin{lettered}
\item The new standardised approach eliminates mandatory modelling for regulatory
capital. Banks now focus modelling on the Internal Capital Adequacy Assessment
Process (ICAAP) under Pillar 2.
\item ICAAP assesses capital sufficiency considering the bank's risk profile and
future plans.
\end{lettered}

Operational risk ICAAP involves scenario analysis:

\begin{checks}
\item Identifying potential operational risk scenarios based on past experiences
and controls.
\item Quantifying potential losses and probabilities for each scenario.
\item Aggregating scenarios to determine the required economic capital.
\end{checks}

% ---------------------------------------------------------------------
\lohead{43}{c}{Describe the differences among key risk indicators (KRIs), key
performance indicators (KPIs), and key control indicators (KCIs).}

\noindent\begin{minipage}{\textwidth}
\renewcommand{\arraystretch}{1.35}
\begin{tabularx}{\textwidth}{@{}L{0.22\textwidth}L{0.25\textwidth}YL{0.22\textwidth}@{}}
\rowcolor{navy} \thd{Indicator} & \thd{What it measures} & \thd{Example} & \thd{Focus}\\
{\tabfont \textbf{1) KRI} \newline (Key Risk Indicator)} &
{\tabfont Risk exposure --- likelihood and impact of potential events} &
{\tabfont Increase in transaction errors; increase in data sensitivity} &
{\tabfont Future risk}\\
\rowcolor{paleblue}
{\tabfont \textbf{2) KPI} \newline (Key Performance Indicator)} &
{\tabfont Performance --- how effectively the bank operates} &
{\tabfont a) Customer complaints \newline b) IT system downtime} &
{\tabfont Current performance}\\
{\tabfont \textbf{3) KCI} \newline (Key Control Indicator)} &
{\tabfont Control effectiveness --- how well controls work} &
{\tabfont Old business continuity plans; uncorrected data entry errors} &
{\tabfont Control effectiveness}\\
\end{tabularx}
\end{minipage}
\orfigcap{KRI, KPI and KCI: three indicators, three time horizons.}

% ---------------------------------------------------------------------
\lohead{43}{g}{Describe and explain the steps to ensure a strong level of
operational resilience, and to test the operational resilience of important
business services.}

Banks need resilience plans to handle disruptions to critical services, not all
processes. This involves identifying those services, figuring out how bad an
outage can be tolerated, and then planning for how to recover.

\begin{orkeybox}
Key steps include identifying critical services, setting tolerances for outages,
mapping dependencies, testing with realistic scenarios, learning from tests and
making adjustments, and having communication plans.
\end{orkeybox}
